\section{The org chart in your schema}
\label{sec:ch01-org-chart}

\index{Conway's law}%
In 1968 Melvin Conway published an observation that has outlived almost
everything else written about software that decade: organisations which design
systems are constrained to produce designs which are copies of the
communication structures of those organisations~\autocite{conway1968committees}.
He was writing about committees and contract negotiations, not APIs, and the
paper is worth reading partly because it is so unlike the slogan it turned
into.

The same paper contains the arithmetic that explains the queue at the end of
section~\ref{sec:ch01-promise}, and I had somehow never noticed it until I went
back to the original. Conway points out that the number of possible
communication paths in an organisation is roughly half the square of the number
of people in it, and draws the conclusion himself: \enquote{Even in a moderately
small organization it becomes necessary to restrict communication in order that
people can get some `work' done.} Restricting communication is what a review
process is. It is not bureaucracy arriving to spoil things; it is the only way a
group above a certain size can ship anything at all.

A GraphQL schema is an unusually literal demonstration. The schema is the
interface, and it is a single document. If six teams contribute to one
document, then those six teams need a communication structure capable of
reaching agreement about that document, and they will build one whether anybody
plans it or not.

You can usually see what they built. Look for a schema working group with a
standing meeting, a naming conventions page that is longer than it should be, a
rule that schema changes need approval from two reviewers outside the
contributing team, an RFC template for adding a type. None of this is
dysfunction. It is a reasonable organisation noticing that a shared artifact
needs shared governance and supplying some. I have helped build these processes
and would build them again.

Both of those descriptions are drawn from life. Netflix ran a schema working
group that met weekly, with a schema architect whose job was to understand the
entire surface area of the graph~\autocite{shin2020netflixtalk}, and they were
straightforward about the trade: \enquote{while reviews add overhead to the
product development process}, they judged that protecting the quality of the
model would cost less later~\autocite{shikhare2020netflix2}. PayPal, once
GraphQL had spread widely enough inside the company, established a standards
body, and reported the bill just as plainly: adoption of the one-graph approach
\enquote{has been slow}, because teams \enquote{have to change a lot of
behavior}, which adds \enquote{process and time to
deliver}~\autocite{kapoor2021paypal}.

They also work. That is the part usually left out of the argument for
federation, so let me put it in early: process solutions to schema coupling are
cheap, effective, and correct for far longer than vendors suggest. If your
complaint is inconsistent field naming, you do not have a distribution problem.
You have a style guide problem, and adopting federation to solve it means paying
for a router and a schema registry to avoid writing a document.
Chapter~\ref{ch:do-you-actually-need-federation} returns to this with the
alternatives worked out properly, because the honest
answer for a lot of teams is that better governance of one graph beats a
federated graph they are not ready to operate.

\index{schema!governance}%
What governance cannot fix is the second arrow in
figure~\ref{fig:ch01-one-schema-many-teams}. You can make review faster. You
can make ownership unambiguous with a \code{CODEOWNERS} file. You cannot make
two teams ship the same binary at different times. Deploy coupling is not a
communication problem, so no amount of communication resolves it, and it is the
one failure mode that gets monotonically worse as the organisation grows.

There is a related trap worth naming, because it is where the coupling starts to
distort the schema itself. When adding a field to an existing type is expensive,
people stop adding fields to existing types. They add
\code{productPricingDetails} alongside \code{product} because a new root field
touches nobody else's code and needs one reviewer. The schema slowly grows a
layer of near-duplicate types whose only reason for existing is that they were
cheaper to merge than the correct change would have been. A year later the graph
is genuinely confusing to clients, and the cause was not incompetence. It was
that the org chart made the right change expensive and the wrong change free.

\index{team autonomy}%
Turned around, this is the actual argument for federation, and it is worth
stating in Conway's terms rather than in terms of performance or scale, neither
of which federation reliably improves. If you want teams that can decide and
ship independently, you need interfaces they can change independently. A single
schema gives you one interface with many owners. Federation gives you many
interfaces, each with one owner, plus machinery that presents them to clients as
though they were still one graph. That machinery is the cost. The independence
is what you are buying.

Whether the trade is worth it depends on how many teams you have and how often
they collide, which is a question about your organisation that no benchmark will
answer. Section~\ref{sec:ch01-do-you-need-this} gives you a way to interrogate
it, and the exercise at the end of the chapter turns it on your own graph.
